\section{Experiments}

\subsection{Experimental protocol}

All main experiments are performed on 1D periodic reaction-diffusion systems. We evaluate three representative cases:
\begin{itemize}[leftmargin=1.5em]
\item \textbf{Case A:} a matched-library reference case with
\begin{equation*}
D(u) = 0.01 + 0.05u, \qquad R(u) = u(1-u).
\end{equation*}
\item \textbf{Case B:} a second matched-library case with
\begin{equation*}
D(u) = 0.01 + 0.03u^2, \qquad R(u) = u - 1.5u^2 + 0.5u^3.
\end{equation*}
\item \textbf{Case Exp:} a function-class-mismatch stress test with
\begin{equation*}
D(u) = 0.01 + 0.035(1-e^{-2.5u}), \qquad R(u) = u(1.1e^{-1.4u} - 0.22).
\end{equation*}
\end{itemize}

We compare four method families:
\begin{enumerate}[leftmargin=1.5em]
\item strong polynomial baseline,
\item weak polynomial baseline,
\item neural numerical surrogate,
\item neural surrogate followed by symbolic compression.
\end{enumerate}
Unless otherwise stated, the neural benchmarks use $8$ training trajectories, $4$-mode random Fourier initial conditions, and the simulation grid $n_x=64$ with integrator step size $\Delta t=10^{-4}$ up to final time $T=0.1$, with snapshots saved every $10$ steps. Models are trained for $150$ epochs. A ``5\% noise'' setting means additive Gaussian perturbations with standard deviation $0.05\,\mathrm{std}(u)$. A ``sparse $x+t$'' setting means spatial stride $2$ and temporal stride $2$. Each reported value is averaged across three random seeds.

\subsection{Baseline map: matched library versus mismatch}

Table~\ref{tab:baseline_map} summarizes representative settings. The main qualitative pattern is immediate. In the matched-library regime, strong and weak polynomial baselines are extremely competitive and behave like correctly specified reference estimators. In the mismatch regime, the neural surrogate becomes relevant as a flexible numerical approximation layer, but its advantage is metric-dependent and does not automatically translate into better rollout.

\begin{table}[t]
\centering
\small
\resizebox{\textwidth}{!}{%
\begin{tabular}{lllccc}
\toprule
Case & Setting & Method & ErrD & ErrR & Unseen \\
\midrule
Case A & clean & strong\_poly & $4.098\mathrm{e}{-04} \pm 1.144\mathrm{e}{-04}$ & $2.113\mathrm{e}{-03} \pm 8.855\mathrm{e}{-04}$ & $1.804\mathrm{e}{-05} \pm 5.128\mathrm{e}{-06}$ \\
Case A & clean & weak\_poly & $8.151\mathrm{e}{-03} \pm 1.047\mathrm{e}{-03}$ & $3.090\mathrm{e}{-02} \pm 9.034\mathrm{e}{-03}$ & $8.307\mathrm{e}{-04} \pm 1.348\mathrm{e}{-04}$ \\
Case A & clean & neural & $6.926\mathrm{e}{-02} \pm 4.007\mathrm{e}{-03}$ & $3.128\mathrm{e}{-01} \pm 1.892\mathrm{e}{-02}$ & $4.828\mathrm{e}{-03} \pm 3.137\mathrm{e}{-04}$ \\
Case A & clean & neural+symbolic & $6.853\mathrm{e}{-02} \pm 4.266\mathrm{e}{-03}$ & $3.127\mathrm{e}{-01} \pm 1.894\mathrm{e}{-02}$ & $4.701\mathrm{e}{-03} \pm 3.516\mathrm{e}{-04}$ \\
\midrule
Case B & clean & strong\_poly & $4.018\mathrm{e}{-04} \pm 1.088\mathrm{e}{-04}$ & $2.597\mathrm{e}{-03} \pm 1.092\mathrm{e}{-03}$ & $1.651\mathrm{e}{-05} \pm 4.630\mathrm{e}{-06}$ \\
Case B & clean & weak\_poly & $8.030\mathrm{e}{-03} \pm 6.990\mathrm{e}{-04}$ & $3.701\mathrm{e}{-02} \pm 8.150\mathrm{e}{-03}$ & $9.556\mathrm{e}{-04} \pm 1.252\mathrm{e}{-04}$ \\
Case B & clean & neural & $1.003\mathrm{e}{-01} \pm 5.687\mathrm{e}{-03}$ & $2.047\mathrm{e}{-01} \pm 1.727\mathrm{e}{-02}$ & $3.949\mathrm{e}{-03} \pm 3.374\mathrm{e}{-04}$ \\
Case B & clean & neural+symbolic & $1.001\mathrm{e}{-01} \pm 5.650\mathrm{e}{-03}$ & $2.036\mathrm{e}{-01} \pm 1.756\mathrm{e}{-02}$ & $3.727\mathrm{e}{-03} \pm 3.585\mathrm{e}{-04}$ \\
\midrule
Case Exp & clean & strong\_poly & $7.664\mathrm{e}{-02} \pm 2.370\mathrm{e}{-03}$ & $3.953\mathrm{e}{-01} \pm 1.225\mathrm{e}{-01}$ & $2.383\mathrm{e}{-03} \pm 2.411\mathrm{e}{-04}$ \\
Case Exp & clean & weak\_poly & $6.395\mathrm{e}{-02} \pm 2.375\mathrm{e}{-03}$ & $1.036\mathrm{e}{-01} \pm 4.733\mathrm{e}{-02}$ & $2.137\mathrm{e}{-03} \pm 1.435\mathrm{e}{-04}$ \\
Case Exp & clean & neural & $4.726\mathrm{e}{-02} \pm 3.732\mathrm{e}{-04}$ & $4.698\mathrm{e}{-01} \pm 8.209\mathrm{e}{-03}$ & $3.162\mathrm{e}{-03} \pm 9.275\mathrm{e}{-05}$ \\
Case Exp & clean & neural+symbolic & $4.710\mathrm{e}{-02} \pm 3.877\mathrm{e}{-04}$ & $4.698\mathrm{e}{-01} \pm 8.209\mathrm{e}{-03}$ & $3.222\mathrm{e}{-03} \pm 8.894\mathrm{e}{-05}$ \\
\bottomrule
\end{tabular}
}
\caption{Representative baseline comparison. In matched-library settings, polynomial baselines are exceptionally strong. In mismatch settings, the neural surrogate serves as a flexible numerical constitutive approximator, while the symbolic stage largely preserves its behavior.}
\label{tab:baseline_map}
\end{table}

\subsection{Excitation coverage and identifiability}

Table~\ref{tab:excitation} makes the identifiability issue explicit. For computational efficiency, this excitation benchmark uses a lighter protocol with $n_x=48$, $T=0.05$, snapshots saved every $5$ steps, and $6$ training trajectories. We compare low-excitation and high-excitation training sets by shrinking or expanding the amplitude range of the random initial conditions. The low-excitation regime occupies only a narrow state interval and produces much weaker gradient and diffusion-energy statistics. In all cases, this regime attains a \emph{smaller} weak loss and even a substantially smaller short-horizon unseen rollout error, yet the recovered closures are substantially worse. Because the signal energy is also much smaller in the low-excitation regime, these weak-loss values should be read as absolute residual magnitudes rather than normalized cross-dataset scores. The conclusion is the one predicted by Proposition~\ref{prop:non_identifiability}: low residual is not enough when the observation operator is weakly excited.

\begin{table}[t]
\centering
\small
\resizebox{\textwidth}{!}{%
\begin{tabular}{llcccccc}
\toprule
Case & Regime & Bin coverage & Weak diffusion energy & Weak loss & ErrD & ErrR & Unseen \\
\midrule
Case A & low excitation & $1.250\mathrm{e}{-01} \pm 0.000\mathrm{e}{+00}$ & $1.091\mathrm{e}{+00} \pm 1.270\mathrm{e}{-01}$ & $1.879\mathrm{e}{-05} \pm 3.018\mathrm{e}{-06}$ & $1.280\mathrm{e}{-01} \pm 9.278\mathrm{e}{-03}$ & $1.289\mathrm{e}{+00} \pm 3.130\mathrm{e}{-02}$ & $5.333\mathrm{e}{-04} \pm 8.114\mathrm{e}{-05}$ \\
Case A & high excitation & $1.000\mathrm{e}{+00} \pm 0.000\mathrm{e}{+00}$ & $1.158\mathrm{e}{+02} \pm 1.368\mathrm{e}{+01}$ & $9.471\mathrm{e}{-04} \pm 1.602\mathrm{e}{-04}$ & $5.488\mathrm{e}{-02} \pm 1.125\mathrm{e}{-02}$ & $3.673\mathrm{e}{-01} \pm 2.049\mathrm{e}{-02}$ & $3.879\mathrm{e}{-03} \pm 3.066\mathrm{e}{-04}$ \\
\midrule
Case B & low excitation & $1.250\mathrm{e}{-01} \pm 0.000\mathrm{e}{+00}$ & $1.822\mathrm{e}{+00} \pm 2.590\mathrm{e}{-01}$ & $1.616\mathrm{e}{-05} \pm 2.869\mathrm{e}{-06}$ & $2.984\mathrm{e}{-01} \pm 1.291\mathrm{e}{-02}$ & $1.031\mathrm{e}{+00} \pm 8.982\mathrm{e}{-02}$ & $4.744\mathrm{e}{-04} \pm 8.218\mathrm{e}{-05}$ \\
Case B & high excitation & $1.000\mathrm{e}{+00} \pm 0.000\mathrm{e}{+00}$ & $2.046\mathrm{e}{+02} \pm 2.852\mathrm{e}{+01}$ & $3.691\mathrm{e}{-04} \pm 5.804\mathrm{e}{-05}$ & $8.566\mathrm{e}{-02} \pm 4.676\mathrm{e}{-03}$ & $2.829\mathrm{e}{-01} \pm 3.448\mathrm{e}{-02}$ & $3.310\mathrm{e}{-03} \pm 1.639\mathrm{e}{-04}$ \\
\midrule
Case Exp & low excitation & $1.250\mathrm{e}{-01} \pm 0.000\mathrm{e}{+00}$ & $2.411\mathrm{e}{+00} \pm 2.782\mathrm{e}{-01}$ & $1.259\mathrm{e}{-05} \pm 1.876\mathrm{e}{-06}$ & $1.706\mathrm{e}{-01} \pm 1.758\mathrm{e}{-02}$ & $8.154\mathrm{e}{-01} \pm 1.803\mathrm{e}{-02}$ & $2.743\mathrm{e}{-04} \pm 1.822\mathrm{e}{-05}$ \\
Case Exp & high excitation & $8.438\mathrm{e}{-01} \pm 0.000\mathrm{e}{+00}$ & $2.185\mathrm{e}{+02} \pm 2.653\mathrm{e}{+01}$ & $3.275\mathrm{e}{-04} \pm 5.214\mathrm{e}{-05}$ & $4.624\mathrm{e}{-02} \pm 1.170\mathrm{e}{-03}$ & $3.912\mathrm{e}{-01} \pm 2.216\mathrm{e}{-02}$ & $2.624\mathrm{e}{-03} \pm 2.115\mathrm{e}{-05}$ \\
\bottomrule
\end{tabular}
}
\caption{Excitation benchmark with three seeds. Narrow low-excitation training sets achieve very small weak losses but much worse closure recovery. The observed trajectory fit alone is therefore not a reliable identifiability diagnostic.}
\label{tab:excitation}
\end{table}

\subsection{Symbolic compression preserves surrogate behavior}

The core symbolic benchmark is summarized in Table~\ref{tab:symbolic_summary} and Figure~\ref{fig:error_propagation}. Across all settings, the symbolic surrogate error remains much smaller than the neural true error. This means the symbolic stage is faithfully compressing the learned constitutive curves rather than inventing a different law. At the same time, the ``noise 5\%'' rows show that weak-form training alone does not make Stage~1 noise-robust: the neural true closure errors deteriorate sharply under observation noise, while the symbolic stage continues to inherit that constitutive bias rather than repair it.

To quantify bias propagation, we define the bias inheritance ratios
\begin{equation}
\mathrm{BIR}_D = \frac{\mathrm{ErrD}^{\mathrm{sym,true}}}{\mathrm{ErrD}^{\mathrm{neural,true}}},
\qquad
\mathrm{BIR}_R = \frac{\mathrm{ErrR}^{\mathrm{sym,true}}}{\mathrm{ErrR}^{\mathrm{neural,true}}}.
\end{equation}
Values near one mean that symbolic compression preserves the constitutive bias already present in the neural surrogate. Values modestly below one indicate mild smoothing, not a systematic correction mechanism.

\begin{table}[t]
\centering
\small
\resizebox{\textwidth}{!}{%
\begin{tabular}{llcccccc}
\toprule
Case & Setting & Neural ErrD & Neural ErrR & Sym surrogate ErrD & Sym surrogate ErrR & $\mathrm{BIR}_D$ & $\mathrm{BIR}_R$ \\
\midrule
Case A & clean & $6.926\mathrm{e}{-02}$ & $3.128\mathrm{e}{-01}$ & $4.628\mathrm{e}{-03}$ & $6.453\mathrm{e}{-03}$ & $9.893\mathrm{e}{-01}$ & $9.999\mathrm{e}{-01}$ \\
Case A & noise 5\% & $8.672\mathrm{e}{-01}$ & $2.361\mathrm{e}{+00}$ & $4.910\mathrm{e}{-04}$ & $4.212\mathrm{e}{-02}$ & $1.000\mathrm{e}{+00}$ & $1.002\mathrm{e}{+00}$ \\
Case A & sparse x+t & $5.046\mathrm{e}{-02}$ & $3.786\mathrm{e}{-01}$ & $1.518\mathrm{e}{-02}$ & $3.676\mathrm{e}{-03}$ & $9.239\mathrm{e}{-01}$ & $1.000\mathrm{e}{+00}$ \\
\midrule
Case B & clean & $1.003\mathrm{e}{-01}$ & $2.047\mathrm{e}{-01}$ & $4.431\mathrm{e}{-03}$ & $2.370\mathrm{e}{-02}$ & $9.988\mathrm{e}{-01}$ & $9.942\mathrm{e}{-01}$ \\
Case B & noise 5\% & $8.470\mathrm{e}{-01}$ & $3.248\mathrm{e}{+00}$ & $5.387\mathrm{e}{-04}$ & $7.419\mathrm{e}{-02}$ & $1.000\mathrm{e}{+00}$ & $1.007\mathrm{e}{+00}$ \\
Case B & sparse x+t & $8.184\mathrm{e}{-02}$ & $2.728\mathrm{e}{-01}$ & $5.720\mathrm{e}{-03}$ & $1.937\mathrm{e}{-02}$ & $9.969\mathrm{e}{-01}$ & $9.981\mathrm{e}{-01}$ \\
\midrule
Case Exp & clean & $4.726\mathrm{e}{-02}$ & $4.698\mathrm{e}{-01}$ & $1.872\mathrm{e}{-03}$ & $1.664\mathrm{e}{-04}$ & $9.965\mathrm{e}{-01}$ & $1.000\mathrm{e}{+00}$ \\
Case Exp & noise 5\% & $8.377\mathrm{e}{-01}$ & $4.580\mathrm{e}{+00}$ & $5.916\mathrm{e}{-04}$ & $2.002\mathrm{e}{-02}$ & $1.000\mathrm{e}{+00}$ & $9.995\mathrm{e}{-01}$ \\
Case Exp & sparse x+t & $7.075\mathrm{e}{-02}$ & $5.442\mathrm{e}{-01}$ & $1.558\mathrm{e}{-03}$ & $1.595\mathrm{e}{-04}$ & $9.990\mathrm{e}{-01}$ & $1.000\mathrm{e}{+00}$ \\
\bottomrule
\end{tabular}
}
\caption{Representative symbolic compression summary. The symbolic stage stays close to the neural surrogate and the bias inheritance ratios remain near one, especially on the mismatch stress test.}
\label{tab:symbolic_summary}
\end{table}

\begin{figure}[t]
\centering
\includegraphics[width=\textwidth]{figures/symbolic_error_propagation.png}
\caption{Neural-to-symbolic error propagation. Stage~2 compression error remains small, while Stage~3 true symbolic error tracks the Stage~1 neural true error. This is the empirical signature of bias preservation.}
\label{fig:error_propagation}
\end{figure}

\subsection{Cross-resolution anti-inverse-crime check}

Table~\ref{tab:cross_resolution} evaluates the learned neural closure on a validation solver with a different grid and saved time step from the data-generation solver. The same-grid unseen rollout and the cross-grid rollout remain nearly identical in all four observation settings. This does not prove full discretization independence, but it does show that the main conclusions are not a trivial artifact of reusing exactly the same numerical configuration at training and validation time.

\begin{table}[t]
\centering
\small
\resizebox{\textwidth}{!}{%
\begin{tabular}{lccccc}
\toprule
Observed stride $(x,t)$ & ErrD & ErrR & Weak loss & Unseen & Cross-grid \\
\midrule
$(1,1)$ & $4.566\mathrm{e}{-02} \pm 9.885\mathrm{e}{-04}$ & $4.064\mathrm{e}{-01} \pm 1.157\mathrm{e}{-02}$ & $2.752\mathrm{e}{-04} \pm 1.139\mathrm{e}{-05}$ & $2.595\mathrm{e}{-03} \pm 4.531\mathrm{e}{-05}$ & $2.580\mathrm{e}{-03} \pm 2.881\mathrm{e}{-05}$ \\
$(2,1)$ & $7.162\mathrm{e}{-02} \pm 1.272\mathrm{e}{-03}$ & $4.841\mathrm{e}{-01} \pm 1.363\mathrm{e}{-02}$ & $2.135\mathrm{e}{-03} \pm 3.685\mathrm{e}{-04}$ & $3.989\mathrm{e}{-03} \pm 1.468\mathrm{e}{-04}$ & $3.975\mathrm{e}{-03} \pm 1.446\mathrm{e}{-04}$ \\
$(1,2)$ & $4.558\mathrm{e}{-02} \pm 9.150\mathrm{e}{-04}$ & $4.092\mathrm{e}{-01} \pm 1.130\mathrm{e}{-02}$ & $2.733\mathrm{e}{-04} \pm 1.018\mathrm{e}{-05}$ & $2.596\mathrm{e}{-03} \pm 4.424\mathrm{e}{-05}$ & $2.581\mathrm{e}{-03} \pm 2.824\mathrm{e}{-05}$ \\
$(2,2)$ & $7.060\mathrm{e}{-02} \pm 1.778\mathrm{e}{-03}$ & $4.823\mathrm{e}{-01} \pm 1.309\mathrm{e}{-02}$ & $1.773\mathrm{e}{-03} \pm 3.379\mathrm{e}{-04}$ & $3.950\mathrm{e}{-03} \pm 1.072\mathrm{e}{-04}$ & $3.936\mathrm{e}{-03} \pm 1.046\mathrm{e}{-04}$ \\
\bottomrule
\end{tabular}
}
\caption{Cross-resolution anti-inverse-crime benchmark on Case~Exp. This benchmark uses a separately generated dataset with a different grid and snapshot cadence from the default protocol. Training data are generated on a fine solver ($n_x=64$, saved $\Delta t=5\times 10^{-4}$), while validation uses a different grid and saved time step ($n_x=48$, saved $\Delta t=6\times 10^{-4}$).}
\label{tab:cross_resolution}
\end{table}

\subsection{Rollout preservation and bias inheritance}

Figure~\ref{fig:rollout_bir} shows the two most important dynamical consequences of symbolic compression. First, the symbolic rollout remains nearly identical to the neural rollout across clean, noisy, and sparse settings. Second, the bias inheritance ratios for both diffusion and reaction remain near one. Together these results support the claim that symbolic compression is not correcting the constitutive law; it is distilling the numerical surrogate into a more compact representation while preserving its existing bias.

\begin{figure}[t]
\centering
\includegraphics[width=\textwidth]{figures/symbolic_rollout_and_bir.png}
\caption{Rollout preservation and bias inheritance. Top: symbolic unseen rollout closely follows neural unseen rollout. Bottom: the bias inheritance ratios remain near one, showing that symbolic compression preserves rather than repairs Stage~1 constitutive bias.}
\label{fig:rollout_bir}
\end{figure}
